% Generated by scripts/assemble.py; corrections belong in source/corrections.json.
\examhead{165}{四川大学 2024 年数学分析试卷(回忆版)}{157}
% p158-b01; source PDF page 158
\par\Needspace{5\baselineskip}
\blockanchor{p158-b01}
\majorquestion{一}计算以下极限.\par
% p158-b02; source PDF page 158
\blockanchor{p158-b02}
\question{1}$\lim_{n\to\infty}\sqrt[n]{1+\sqrt[n]{2+\cdots+\sqrt[n]{n}}}$ .\par
% p158-b03; source PDF page 158
\blockanchor{p158-b03}
\question{2}$\lim_{x\to0}\left(\frac{x}{\sin x}-\cos x+1\right)\frac{1}{x}.$\par
% p158-b04; source PDF page 158
\blockanchor{p158-b04}
\question{3}$\lim_{n\to\infty}\frac{\ln2+\ln3+\cdots+\ln n}{n^{2}\ln n}.$\par
% p158-b05; source PDF page 158
\blockanchor{p158-b05}
\question{4}$\lim_{x\to 0^{+}}\frac{x}{\cos x}\int_{x^{-1}}^{x^{-2}}\left(1 + \frac{1}{2t}\right)^{t}\sin \frac{1}{\sqrt{t}}\mathrm{d}t.$\par
% p158-b06; source PDF page 158
\par\Needspace{5\baselineskip}
\blockanchor{p158-b06}
\majorquestion{二}计算下列积分.\par
% p158-b07; source PDF page 158
\blockanchor{p158-b07}
\question{1}$\int \frac{1 + x}{1 - x^4} \, \mathrm{d}x.$\par
% p158-b08; source PDF page 158
\blockanchor{p158-b08}
\question{2}$\oint_{L} \frac{y}{(2 - x)^{2} + y^{2}} \, \mathrm{d}y + \frac{x}{(2 - x)^{2} + y^{2}} \, \mathrm{d}y,$ 其中 $L$ 是不经过 (2,0) 点的封闭曲线, 方向取逆时针.\par
% p158-b09; source PDF page 158
\blockanchor{p158-b09}
\question{3}$\iint_{\Sigma} \frac{2}{x} \mathrm{d}y \mathrm{d}z + \mathrm{d}z \mathrm{d}x + \frac{1}{z} \mathrm{d}x \mathrm{d}y,$ 其中 $\Sigma$ 是 $x^{2} + y^{2} + z^{2} = 1$ 的表面, 方向取外侧.\par
% p158-b10; source PDF page 158
\blockanchor{p158-b10}
\question{4}$\int_0^{+\infty}\frac{e^{-ax} - e^{-bx}}{x}\mathrm{d}x.$\par
% p158-b11; source PDF page 158
\par\Needspace{5\baselineskip}
\blockanchor{p158-b11}
\majorquestion{三}证明区间上有第一类间断点的函数一定没有原函数.\par
% p158-b12; source PDF page 158
\par\Needspace{5\baselineskip}
\blockanchor{p158-b12}
\majorquestion{四}设 $f(x)$ 在 (0,1) 上有界且可导, $\lim_{x\to 0^{+}}f(x)$ 不存在, 证明一定存在数列 $\{x_n\}$, 当 $\lim_{n\to \infty}x_n = 0$ 时一定有 $\lim_{n\to \infty}f(x_n) = 0$.\par
% p158-b13; source PDF page 158
\par\Needspace{5\baselineskip}
\blockanchor{p158-b13}
\majorquestion{五}设 $F(x, y) = 0$ 是平面上一条八字型光滑曲线, 说明曲线上满足 $\frac{\partial F}{\partial x} = \frac{\partial F}{\partial y} = 0$ 的点至少有几个兵说明理由.\par
% p158-b14; source PDF page 158
\par\Needspace{5\baselineskip}
\blockanchor{p158-b14}
\majorquestion{六}连续可微函数 $z = z(x, y)$ 由 $F(xz - y, x - yz) = 0$ 唯一确定.\par
% p158-b15; source PDF page 158
\blockanchor{p158-b15}
(1) 求  $\frac{\partial z}{\partial x}, \frac{\partial z}{\partial y}.$\par
% p158-b16; source PDF page 158
\blockanchor{p158-b16}
(2) 证明: $\oint_{L} (xy^{2} + 2yz) \, \mathrm{d}y - (2xz + yz^{2}) \, \mathrm{d}x = 2\pi$, $L$ 为 $x^{2} + y^{2} = 1$, 方向取逆时针.\par
% p158-b17; source PDF page 158
\par\Needspace{5\baselineskip}
\blockanchor{p158-b17}
\majorquestion{七}讨论反常积分 $\int_0^{+\infty} \frac{\ln x}{x} \sin x \, \mathrm{d}x$ 的收敛性(包括条件收敛和绝对收敛).\par
% p158-b18; source PDF page 158
\par\Needspace{5\baselineskip}
\blockanchor{p158-b18}
\majorquestion{八}证明 $f(x, y, z) = \frac{1}{1 - xyz}$ 在 $[0, 1) \times [0, 1) \times [0, 1)$ 上非一致连续.\par
% p158-b19; source PDF page 158
\par\Needspace{5\baselineskip}
\blockanchor{p158-b19}
\majorquestion{九}设 $x < f_{1}(x) < \sqrt{x}$, $f_{n}(x) = \sqrt{f_{n-1}(x)}$.\par
% p158-b20; source PDF page 158
\blockanchor{p158-b20}
(1) 对任意的 $x_0 \in (0,1)$, 证明: $\{f_n(x_0)\}$ 是单调有界数列.\par
% p158-b21; source PDF page 158
\blockanchor{p158-b21}
(2) 证明: $\{f_n(x)\}$ 在 (0,1) 上一致收敛.\par
% p158-b22; source PDF page 158
\par\Needspace{5\baselineskip}
\blockanchor{p158-b22}
\majorquestion{十}想不起来了.\par
